\section{Discussion}
\label{sec:discussion}

\paragraph{Difficulty definitions are not interchangeable.}
Task and Expert difficulty are strongly related ($\rho=0.842$) because the
deterministic expert converts longer approach and transport geometry into
longer executions. Learner difficulty instead measures normalized motion error
of a particular fitted MLP family. It correlates negatively with duration
($\rho=-0.542$), and 43 of 80 Learner-Hard trajectories are Task-Easy. Grasp
and release have the largest phasewise errors, but the score is a per-step
average; these facts make trajectory length and manipulation-phase composition
plausible contributors to ranking disagreement, not established causes.

\paragraph{Staged curricula combine ordering and reweighting.}
Forward schedules expose the first group to about 60\% of all sampled
transitions and the last to only 11--12\%, whereas Reverse Learner assigns
60.1\% to Learner-Hard. The comparison therefore does not isolate ordering
from cumulative exposure. The random-ranking control preserves this staged
mechanism but changes membership: its 69.0\% development success lies below
original Reverse Learner (84.7\%) and Uniform (75.7\%). One corruption per
level, a non-monotonic response, and three seeds prevent interpreting this as
a clean causal effect of ranking quality.

\paragraph{Reverse Learner is uncertain and setting-dependent.}
Its pick-and-place Fresh-ID mean is highest (81.5\% versus 69.3\% for
Uniform), but paired seed differences are $+9.5$, $+28.0$, and $-1.0$
percentage points. Development AUCs are nearly equal (0.787 versus 0.784), and
some intermediate Uniform checkpoints perform strongly; the final mean does
not imply faster convergence. In exploratory stacking, every method is near
the floor and no curriculum exceeds Uniform. Stacking also changes
representation and applies precision-phase emphasis to every schedule, so the
cross-task contrast cannot be attributed solely to task identity.

\paragraph{Closed-loop outcomes matter.}
Final validation losses occupy a narrow range while rollout success varies
widely. One-step imitation metrics therefore remain diagnostics rather than
surrogates for physical completion: errors in grasping, alignment, release, or
retreat alter subsequent states. This gap directly addresses RQ2 and motivates
evaluating curriculum schedules through autonomous closed-loop success.
